\documentclass[11pt]{article}

\usepackage[margin=1in]{geometry}
\usepackage{amsmath,amssymb,amsthm}
\usepackage{hyperref}
\usepackage{microtype}

\title{\textbf{Infinity Wall: VCDI Instantiation and Directed-Limit Rigidity}}
\author{Inacio F.\ Vasquez \\ Independent Researcher}
\date{}

\newtheorem{theorem}{Theorem}
\newtheorem{lemma}[theorem]{Lemma}
\newtheorem{proposition}[theorem]{Proposition}
\newtheorem{corollary}[theorem]{Corollary}
\theoremstyle{definition}
\newtheorem{definition}[theorem]{Definition}
\theoremstyle{remark}
\newtheorem{remark}[theorem]{Remark}

\newcommand{\ED}{\mathrm{ED}}
\newcommand{\F}{\mathbb{F}}
\newcommand{\R}{\mathbb{R}}

\begin{document}
\maketitle

\begin{abstract}
We record a canonical instantiation of the Vasquez Capacity--Depth Inequality (VCDI) on expander-like and random-lift graph families, yielding an ``infinity wall'' statement: under URF admissibility with finite per-step transcript capacity, any exhaustion with unbounded entropy demand forces unbounded entropy depth, hence rigidity in the directed limit.
\end{abstract}

\section{Admissible class and capacity}
Fix an admissible refinement class \(\mathcal{A}\) with finite per-step transcript capacity bound \(C_{\max}\in(0,\infty)\).
For each finite instance \(P\in\mathcal{A}\), let \(H(P)\) denote its entropy demand and \(\ED(P)\) its entropy depth.

\section{VCDI (canonical form)}
\begin{theorem}[Vasquez Capacity--Depth Inequality (VCDI)]
\label{thm:vcdi}
For every admissible process \(P\in\mathcal{A}\),
\[
\ED(P)\ \ge\ \frac{H(P)}{C_{\max}}.
\]
\end{theorem}

\section{Concrete instances from graph families}

\subsection{Cycle-rank invariant}
Let \(G\) be a finite connected graph. Write \(Z_1(G;\F_2)\) for the \(\F_2\)-cycle space and define the global invariant
\[
I(G)\ :=\ \dim_{\F_2} Z_1(G;\F_2).
\]
For a connected graph with \(n\) vertices and \(m\) edges, \(I(G)=m-n+1\).

\subsection{Expander / high cycle-rank families}
Let \((G_n)_{n\ge 1}\) be a family of connected \(d\)-regular graphs (\(d\ge 3\)) with \(|V(G_n)|=n\).
Then \(m=\tfrac{dn}{2}\) and hence
\[
I(G_n)\ =\ m-n+1\ =\ \Big(\frac d2-1\Big)n+1\ =\ \Theta(n).
\]

\begin{definition}[Expander-instantiated process]
For each \(n\), let \(P_n\in\mathcal{A}\) be a URF-admissible refinement process whose entropy demand includes the uncertainty of a uniform parity state on \(Z_1(G_n;\F_2)\).
\end{definition}

\begin{lemma}[Linear entropy demand from cycle parity]
\label{lem:linearH}
Under the above definition,
\[
H(P_n)\ \ge\ I(G_n)\ =\ \Theta(n).
\]
\end{lemma}

\begin{corollary}[Linear depth lower bound on \(d\)-regular families]
\label{cor:linearED}
For such \((P_n)\subset\mathcal{A}\),
\[
\ED(P_n)\ \ge\ \frac{H(P_n)}{C_{\max}}\ \ge\ \frac{I(G_n)}{C_{\max}}\ =\ \Theta(n).
\]
\end{corollary}

\subsection{Random lifts}
Let \(B\) be a fixed connected base graph with minimum degree \(\ge 3\). Let \(G_n\) be an \(n\)-lift of \(B\).
Since \(|E(G_n)|=n|E(B)|\) and \(|V(G_n)|=n|V(B)|\),
\[
I(G_n)\ =\ |E(G_n)|-|V(G_n)|+1\ =\ n(|E(B)|-|V(B)|)+1\ =\Theta(n)
\]
provided \(|E(B)|-|V(B)|>0\) (equivalently \(B\) is not a tree).

\begin{corollary}[Random-lift instantiation]
\label{cor:lifts}
For any admissible \(P_n\in\mathcal{A}\) whose demand includes uniform cycle-parity uncertainty on \(Z_1(G_n;\F_2)\),
\[
H(P_n)=\Theta(n),
\qquad
\ED(P_n)\ge \Theta(n).
\]
\end{corollary}

\section{Sharp constant for fixed admissible class}
\begin{proposition}[Sharpness of \(1/C_{\max}\) for a fixed class]
\label{prop:sharp}
Assume \(\mathcal{A}\) is specified so that for every \(P\in\mathcal{A}\) and every step \(t\),
\[
\Delta H_t(P)\le C_{\max}.
\]
Then Theorem~\ref{thm:vcdi} holds, and the factor \(1/C_{\max}\) is sharp in the sense that equality is achieved by processes satisfying \(\Delta H_t(P)=C_{\max}\) at every step until termination.
\end{proposition}

\section{Global invariant separator}
\begin{proposition}[Computability separation witness]
\label{prop:sep}
The invariant \(I(G)=\dim_{\F_2}Z_1(G;\F_2)\) is computable by global aggregation. Any regime that terminates by explicitly computing \(I(G_n)\) is non-URF, in that it violates locality (FO\(^k\)-local refinement with fixed \(k\)) when used as a termination-critical quantity.
\end{proposition}

\section{Directed-limit dichotomy (Infinity Wall)}
Let \((P_n)_{n\ge 1}\subset\mathcal{A}\) be any exhaustion such that \(H(P_n)\uparrow\infty\).

\begin{theorem}[Infinity Wall]
\label{thm:infwall}
If \(C_{\max}<\infty\) and VCDI holds in the form of Theorem~\ref{thm:vcdi}, then
\[
\lim_{n\to\infty}\ED(P_n)\ =\ \infty.
\]
Define a directed limit object \(P_\infty:=\varinjlim P_n\). Then \(\ED(P_\infty)=\infty\); i.e.\ the limit is rigid (non-terminating) within the admissible refinement regime.
\end{theorem}

\begin{proof}
By VCDI, \(\ED(P_n)\ge H(P_n)/C_{\max}\). Since \(H(P_n)\uparrow\infty\) and \(C_{\max}<\infty\), the right-hand side diverges, hence \(\ED(P_n)\to\infty\). The final assertion is the definition of \(\ED(P_\infty)\) as the directed-limit divergence of \(\ED(P_n)\).
\end{proof}

\section{Classification rule}
An ``infinite'' method that succeeds must exit URF by at least one of:
\begin{itemize}
\item capacity violation: \(C_{\max}=\infty\) (unbounded transcript per step),
\item locality violation: reliance on global invariants (e.g.\ \(I(G)\), spectrum, global homology).
\end{itemize}
Otherwise (URF-admissible with \(C_{\max}<\infty\)), Theorem~\ref{thm:infwall} forces rigidity.

\end{document}

